% !TeX root = ../navier-stokes-manuscript.tex
% ============================================================
% 1.2 Historical Background
% ============================================================
\subsection{Historical Background}\label{sub:HistoricalEstimates}

Once a smooth solution exists locally, the whole problem is keeping enough of it under control to restart it past every later time.
The historical estimates matter because they show what the equations can already carry that far, what each kind of control buys, and what still has room to escape.
The first quantity the equations keep under global control is kinetic energy.

\divider{Leray's Energy Estimate}

For a smooth decaying solution, pair the momentum equation with \(u\) and integrate over space.
Nonlinear transport moves kinetic energy through the fluid, but incompressibility makes its total contribution vanish.
Pressure redirects the motion without changing the total energy, so its contribution vanishes as well.
Viscosity is the term that remains, and it removes energy through spatial variation:
\begin{equation}\label{eq:LerayEnergyInequality}
  \lVert u(t)\rVert_{L^2}^{2}
  +2\nu\int_0^t\lVert\nabla u(s)\rVert_{L^2}^{2}\,ds
  \leq
  \lVert u_0\rVert_{L^2}^{2}.
\end{equation}
The first term keeps the kinetic energy bounded at each time.
The integral records the total spatial variation that viscosity dissipates over the interval.

Leray's decisive step was to carry this inequality through a limiting sequence of smooth approximations.
Weak lower semicontinuity preserves the inequality even when the approximating derivatives do not converge strongly.
That was enough to construct a finite-energy weak solution on \(\mathbb R^3\) for all time \parencite{Leray1934}.
In modern norm language,
\[
  u\in L^\infty(0,T;L^2),
  \qquad
  \nabla u\in L^2(0,T;L^2)
\]
on every finite interval.
The solution reaches every finite time in the energy class, but the estimate does not retain the classical derivatives or give uniqueness in three dimensions.
The time interval is global.
The regularity is not.

\divider{The Enstrophy Problem}

The dissipation in Leray's estimate already contains the next quantity we need to understand.
For a divergence-free decaying field,
\[
  \lVert\nabla u\rVert_{L^2}^{2}
  =
  \lVert\omega\rVert_{L^2}^{2},
  \qquad
  \omega:=\nabla\times u.
\]
The squared \(L^2\) size of the vorticity is the enstrophy.
Leray's estimate pays for its total accumulation over time.
It does not cap its height at each individual time.
A finite integral can still contain very tall concentrations on very short intervals.

Taking the curl of the momentum equation shows why this first derivative can grow:
\begin{equation}\label{eq:HistoricalVorticityEquation}
  \partial_t\omega+(u\cdot\nabla)\omega
  =
  (\omega\cdot\nabla)u+\nu\Delta\omega.
\end{equation}
Transport carries the vorticity with the fluid.
Viscosity diffuses it.
The remaining term stretches it.
In two dimensions, vortex stretching is absent and the enstrophy estimate closes.
In three dimensions, the surrounding velocity gradient can elongate a vortex line and increase its intensity while viscosity is trying to smooth it.
The balance becomes
\[
  \frac12\frac{d}{dt}\lVert\omega\rVert_{L^2}^{2}
  +\nu\lVert\nabla\omega\rVert_{L^2}^{2}
  =
  \int_{\mathbb R^3}(\omega\cdot\nabla)u\cdot\omega\,dx.
\]
Viscosity has a sign in this equation.
Vortex stretching does not.
The global energy bound does not compare them strongly enough to keep the enstrophy uniformly bounded in time.

\divider{Higher Sobolev Continuation}

Once a derivative-level quantity can escape, local theory tells us how much control would be enough to keep the smooth solution going.
Suppose the solution is smooth on \([0,T)\) and, for some \(s>5/2\),
\begin{equation}\label{eq:HistoricalSobolevContinuation}
  \sup_{0\leq t<T}\lVert u(t)\rVert_{H^s}<\infty.
\end{equation}
Sobolev embedding then keeps \(\nabla u\) bounded, the differentiated energy estimates stay controlled, and the local theory restarts the solution past \(T\).
The same estimates also prevent a finite endpoint whenever
\[
  \int_0^T\lVert\nabla u(t)\rVert_{L^\infty}\,dt<\infty.
\]
These continuation principles grew out of the classical Hilbert and Sobolev theories and their later scale-sensitive \(L^p\) versions \parencite{FujitaKato1964,Kato1984}.

This changes what a finite endpoint means.
A smooth solution cannot reach \(T\), keep one of these restart conditions finite, and still stop there.
The continuation results tell us exactly what kind of control would be enough.
They do not derive that control from arbitrary smooth initial data for every finite time.

\divider{Caffarelli\textendash Kohn\textendash Nirenberg Partial Regularity}

Without a global continuation bound, the energy-class solution still carries local information.
Caffarelli, Kohn, and Nirenberg moved the energy argument from the whole space into parabolic cylinders.
Around a spacetime point \(z_0=(x_0,t_0)\), write
\[
  Q_r(z_0)
  :=
  B_r(x_0)\times(t_0-r^2,t_0).
\]
Suitable weak solutions satisfy a local energy inequality on these cylinders.
Their epsilon-regularity principle says that sufficiently small scale-invariant velocity and pressure quantities on \(Q_r(z_0)\) force the solution to be regular on a smaller cylinder.

The same statement read from a possible singular point has a sharp consequence.
Every sufficiently small cylinder around that point must retain a definite amount of scale-invariant concentration.
If the concentration fell below the epsilon threshold at one scale, the smaller cylinder would already be regular.
A covering argument then confines the possible singular set to a set of zero one-dimensional parabolic Hausdorff measure \parencite{CaffarelliKohnNirenberg1982}.
That is a severe restriction.
It still allows a point.
One point is enough to break global smoothness.

\divider{Ladyzhenskaya\textendash Prodi\textendash Serrin Criteria}

The local picture makes the next control problem concrete.
A possible singular point has to sustain enough velocity concentration in space for long enough in time.
The Ladyzhenskaya\textendash Prodi\textendash Serrin estimates measure those two demands together.

For \(q>3\), H\"older's inequality, interpolation, and Young's inequality turn the first differentiated energy estimate into
\[
  \frac{d}{dt}\lVert\nabla u\rVert_{L^2}^{2}
  +\nu\lVert\Delta u\rVert_{L^2}^{2}
  \leq
  C_{\nu,q}\lVert u\rVert_{L^q}^{\frac{2q}{q-3}}
  \lVert\nabla u\rVert_{L^2}^{2}.
\]
The coefficient on the right is integrable in time when
\begin{equation}\label{eq:HistoricalLPSCriterion}
  u\in L^p(0,T;L^q(\mathbb R^3)),
  \qquad
  \frac{2}{p}+\frac{3}{q}\leq1,
  \qquad
  q>3.
\end{equation}
Gr\"onwall's inequality then keeps the enstrophy finite, and the continuation theory keeps the solution regular through \(T\).
This is the Ladyzhenskaya\textendash Prodi\textendash Serrin criterion for Leray\textendash Hopf weak solutions \parencite{Prodi1959,Serrin1963,Ladyzhenskaya1967}.

The exponent \(q\) measures how sharply velocity can concentrate in space.
The exponent \(p\) measures how long that concentration can persist.
At equality, the mixed norm is unchanged by the Navier\textendash Stokes scaling; strict inequality assumes stronger control.
The criterion lets spatial intensity trade against temporal duration while still protecting the solution.
Any finite endpoint has to escape every one of these balances.

\divider{The Escauriaza\textendash Seregin\textendash \v{S}ver{\'a}k Endpoint}

The direct Serrin estimate runs out exactly at \(q=3\), because the exponent \(2q/(q-3)\) becomes infinite.
Scaling makes that endpoint impossible to dismiss.
Under the Navier\textendash Stokes rescaling
\begin{equation}\label{eq:HistoricalNavierStokesScaling}
  u_\lambda(x,t)
  :=
  \lambda u(\lambda x,\lambda^2t),
\end{equation}
the spatial \(L^3\) norm is unchanged:
\[
  \lVert u_\lambda(t)\rVert_{L^3}
  =
  \lVert u(\lambda^2t)\rVert_{L^3}.
\]
The critical norm neither shrinks nor grows when a suspected concentration is magnified back to unit scale.

Escauriaza, Seregin, and \v{S}ver{\'a}k proved that a Leray\textendash Hopf weak solution on \(\mathbb R^3\) remains regular through \(T\) whenever
\[
  u\in L^\infty(0,T;L^3(\mathbb R^3))
\]
\parencite{EscauriazaSereginSverak2003}.
The proof assumes a first singular time and magnifies the concentration around it.
Because the \(L^3\) bound survives that magnification, compactness produces a nontrivial ancient solution with the same critical control.
Backward uniqueness for its vorticity then forces the limiting solution to vanish, contradicting the concentration used to produce it.
The endpoint closes through rigidity after rescaling even though the direct Serrin estimate no longer closes there.

Later results made the endpoint demand sharper.
For a smooth finite-energy solution on \(\mathbb R^3\) with finite maximal time \(T_*\), Seregin proved
\begin{equation}\label{eq:SereginLThreeDivergence}
  \lim_{t\uparrow T_*}\lVert u(t)\rVert_{L^3(\mathbb R^3)}
  =\infty
\end{equation}
\parencite{Seregin2012}.
The critical norm cannot rise along a few isolated times and then return below a fixed level.
It must eventually leave every finite level.
Tao later proved that, along a sequence approaching any finite maximal time, the same norm must exceed a positive power of a triple logarithm of the inverse time remaining \parencite{Tao2021Quantitative}.
That rate grows extremely slowly.
It still forces a possible finite endpoint to carry increasing critical concentration as the available time disappears.

\divider{Where the Problem Sits Today}

The scaling used at the endpoint also explains why Leray's global estimate does not supply the critical bound.
At corresponding times,
\[
  \lVert u_\lambda(t)\rVert_{L^3}
  =
  \lVert u(\lambda^2t)\rVert_{L^3},
  \qquad
  \lVert u_\lambda(t)\rVert_{L^2}^{2}
  =
  \lambda^{-1}\lVert u(\lambda^2t)\rVert_{L^2}^{2}.
\]
Magnifying a smaller spatial scale leaves the critical intensity unchanged while the rescaled kinetic energy becomes smaller.
The energy inequality controls how much velocity is present in total.
The critical criteria control how tightly that velocity can concentrate.
The published estimates do not yet convert the first control into the second for arbitrary smooth data.

Nonlinear transport can move activity toward smaller spatial scales, where the gradients become larger and viscosity acts more strongly.
The remaining difficulty is proving that viscous dissipation must control that transfer before any continuation norm can escape.
As a heuristic only, suppose the active frequency doubles from one stage to the next.
The corresponding spatial scale halves, and its natural parabolic time scale quarters.
Those model times form a convergent geometric series, so an unlimited sequence of scale changes can fit inside a finite total time.
The Zeno description refers to that finite sum.
It does not assert that a Navier\textendash Stokes solution follows the model.

A global regularity estimate has to remain effective while the active scale changes.
It has to come from the initial data, the viscosity, and the equations, and it has to keep the solution inside a continuation regime through every finite time.
Publicly established theory has not yet supplied that bound for arbitrary smooth finite-energy data.
